\DocumentMetadata{
  pdfversion=2.0,
  pdfstandard=ua-2,
  lang=en-US,
  tagging=on,
  tagging-setup={math/setup=mathml-SE,tabsorder=structure}
}
\documentclass[11pt,letterpaper]{article}
\usepackage[margin=0.68in,includefoot]{geometry}
\usepackage{newcomputermodern}
\usepackage{amsmath,amscd,mathtools}
\providecommand{\square}{\mdlgwhtsquare}
\usepackage[hidelinks]{hyperref}
\hypersetup{
  pdftitle={Topology PhD Exam, September 1993},
  pdfauthor={University of Florida Department of Mathematics},
  pdfsubject={Graduate mathematics examination},
  pdfdisplaydoctitle=true,
  bookmarksopen=true
}
\setcounter{secnumdepth}{0}
\setlength{\parindent}{0pt}
\setlength{\parskip}{0.28em}
\setlength{\emergencystretch}{3em}
\setlength{\leftmargini}{2.15em}
\setlength{\leftmarginii}{2.65em}
\setlength{\leftmarginiii}{3em}
\renewcommand{\labelenumi}{\textbf{\theenumi.}}
\pagestyle{empty}
\raggedbottom
\allowdisplaybreaks
\newenvironment{examproblems}[1]{%
  \begin{enumerate}%
  \setcounter{enumi}{#1}%
  \setlength{\topsep}{0.35em}%
  \setlength{\partopsep}{0pt}%
  \setlength{\itemsep}{0.48em}%
  \setlength{\parsep}{0.18em}%
}{\end{enumerate}}
\newenvironment{examsubparts}{%
  \begin{enumerate}%
  \setlength{\topsep}{0.18em}%
  \setlength{\partopsep}{0pt}%
  \setlength{\itemsep}{0.16em}%
  \setlength{\parsep}{0.1em}%
}{\end{enumerate}}
\newenvironment{examinstructions}{%
  \par\begingroup\itshape
}{%
  \par\endgroup\vspace{0.18em}
}
\makeatletter
\renewcommand\section{\@startsection{section}{1}{0pt}%
  {-1.25ex plus -.25ex minus -.1ex}%
  {0.55ex plus .1ex}%
  {\normalfont\large\bfseries}}
\renewcommand{\@maketitle}{%
  \newpage\null\vskip -1.2em%
  {\footnotesize\bfseries
    \href{https://www.ufl.edu/}{University of Florida}, \href{https://math.ufl.edu/}{Department of Mathematics}\par}%
  \vskip 0.18em%
  \tagstructbegin{tag=Title}%
    {\LARGE\bfseries\href{https://gma.math.ufl.edu/exams/topology-phd/}{Topology PhD Exam}, September 1993\par}%
  \tagstructend%
  \vskip 0.35em%
  \tagmcbegin{artifact}%
    \hrule height 1pt%
  \tagmcend%
  \vskip 0.55em%
}
\makeatother
\title{Topology PhD Exam, September 1993}
\author{}
\date{}
\begin{document}
\maketitle
\thispagestyle{empty}
\begin{examinstructions}
Answer the following questions without giving proofs.
\end{examinstructions}
\begin{examproblems}{0}
\item
Let \(X \subset Y\).
\begin{examsubparts}
\item[\textnormal{a)}]
What does it mean for~\(X\) to be a retract of~\(Y\)?
\item[\textnormal{b)}]
What does it mean for~\(X\) to be a deformation retract of~\(Y\)?
\end{examsubparts}
\item
Let~\(K\) be a finite simplicial complex. Define \(\chi(K)\), the Euler characteristic of~\(K\).
\item
Define what it means for a topological space to be
\begin{examsubparts}
\item[\textnormal{a)}]
normal
\item[\textnormal{b)}]
connected
\end{examsubparts}
\item
State the Brouwer fixed point theorem.
\end{examproblems}
\begin{examinstructions}
Answer the following questions, giving proofs or counterexamples.
\end{examinstructions}
\begin{examproblems}{4}
\item
Let \(p:X \to Y\) be a~\(d\)-sheeted covering projection, where~\(X\) and~\(Y\) are finite simplicial complexes. Show that \(\chi(X)=d\cdot\chi(Y)\).
\item
Compute the fundamental group of a closed, orientable surface~\(X\) of genus 3.
\item
Let~\(X\) be a closed, orientable surface of genus 3.
\begin{examsubparts}
\item[\textnormal{a)}]
Compute the homology groups of~\(X\).
\item[\textnormal{b)}]
Compute the Euler characteristic of this surface~\(X\).
\end{examsubparts}
\item
Show that every continuous map \(f:S^2 \to S^1\) is homotopic to a constant map.
\item
For \(i=1,2,\ldots\), let \(P_i\) be a 2-dimensional plane in \(\mathbb{R}^3\).
\begin{examsubparts}
\item[\textnormal{a)}]
Show that \(S=\mathbb{R}^3\setminus\bigcup_{i=1}^{\infty}P_i\) is not empty.
\item[\textnormal{b)}]
Must~\(S\) be an open set?
\end{examsubparts}
\item
Let~\(C\) be the Cantor set and let~\(S\) be any compact metric space. Show that there is a continuous map \(f:C \to S\) such that~\(f\) is onto.
\end{examproblems}
\end{document}
